\documentclass[11pt]{article}
\usepackage[margin=1in]{geometry}
\usepackage[T1]{fontenc}
\usepackage{microtype}
\usepackage{hyperref}
\usepackage{xcolor}
\usepackage{enumitem}
\usepackage{booktabs}
\usepackage{tabularx}
\usepackage{ragged2e}

\hypersetup{
  pdftitle={Advisor Email Record - Research Direction},
  pdfauthor={Nicholas Mino},
  pdfsubject={Emails that shaped the research direction, May--July 2026},
  colorlinks=false,
  pdfborder={0 0 0}
}

\definecolor{critical}{RGB}{140,20,20}
\definecolor{superseded}{RGB}{120,70,20}
\definecolor{current}{RGB}{20,90,50}
\definecolor{editorial}{RGB}{60,60,60}

\clubpenalty=10000
\widowpenalty=10000
\displaywidowpenalty=10000

\newcommand{\CriticalSource}[1]{%
  \par\noindent\textcolor{critical}{\textbf{EXTREMELY IMPORTANT SOURCE --- #1}}%
  \par\medskip
}
\newcommand{\Meta}[2]{%
  \par\noindent\textbf{#1:} #2\par\nobreak
}
\newcommand{\Email}[3]{%
  \medskip\noindent\textbf{On #1, #2 wrote:}\par\nobreak
  \begingroup\RaggedRight #3\par\endgroup
}
\newcommand{\Editorial}[1]{%
  \par\medskip\noindent\textcolor{editorial}{\textbf{Editorial note.} #1}\par\medskip
}
\newcommand{\Takeaway}[1]{%
  \par\medskip
  \noindent\fbox{\parbox{\dimexpr\linewidth-2\fboxsep-2\fboxrule}{%
    \RaggedRight\textcolor{editorial}{\textbf{Central evidentiary conclusion.} #1}%
  }}\par
}

\title{Advisor Email Record\\[0.35em]
\large Emails That Shaped the Research Direction\\[0.15em]
\normalsize May--July 2026}
\author{Nicholas Mino\\
Carnegie Mellon University\\[0.4em]
\small Faculty mentor: Woody Zhu\\
\small Graduate mentor: Wenbin Zhou}
\date{Document compiled 26 July 2026\\
\small Source PDF: \texttt{Advisor\_Email\_Record\_Research\_Direction.pdf}}

\begin{document}
\maketitle

\begin{center}
\textit{Primary documentary record. Advisor-authored guidance is distinguished from
student proposals and administrative context.}
\end{center}

\section*{Purpose and Provenance}
This document contains the email threads supplied by Nick Mino that materially
shaped, clarified, or documented the evolution of the research project. The
emails are ordered chronologically. Each entry begins with an annotation
explaining its authority, significance, and relationship to the developing
research direction. Routine scheduling material is omitted unless it
establishes the occurrence or context of an important research meeting.

\section*{How to Read This Record}
\begin{itemize}[leftmargin=*,itemsep=0.35em]
  \item Advisor-authored guidance has the highest evidentiary weight.
  \item Student-authored descriptions document the proposal presented to the
        advisors, but do not by themselves establish advisor approval.
  \item Later advisor formulations may supersede earlier student formulations.
  \item The June~11 email thread establishes the location and occurrence of an
        in-person meeting; substantive content from that meeting comes from
        contemporaneous notes and belongs in the main research packet, not in
        this email record.
\end{itemize}

\section*{Executive Summary}
The email record shows a clear progression. Professor Zhu first directed Nick
toward the group's uncertainty-quantification and decision-focused papers and
connected him with Wenbin Zhou for closer technical guidance. By late June,
Nick proposed studying stabilized selection algorithms through their primitive
operations and recurring stabilization mechanisms. After the June~29 meeting,
Wenbin reported that he had rewritten the problem setup in Overleaf so that it
aligned more closely with what the advisors intended, and instructed Nick to
begin developing methodologies for that formulation. A June~12 administrative
exchange also records Wenbin's recommendation that the project description
explicitly use the term ``post-hoc inference.''

\Takeaway{The current formulation should be grounded in Wenbin's July~2
advisor-authored problem setup and the June~29 discussion, while Nick's June~25
primitive-operation framework should be treated as a student proposal that
motivated that discussion rather than as the final approved formulation.}

\section*{Chronological Timeline}
\begin{center}
\begin{tabularx}{\linewidth}{@{}l >{\RaggedRight\arraybackslash}X
  >{\RaggedRight\arraybackslash}X >{\RaggedRight\arraybackslash}X@{}}
\toprule
\textbf{Date} & \textbf{Authority} & \textbf{Importance} &
  \textbf{Research significance} \\
\midrule
May 11--14 & Student request / advisor direction & High &
  Professor Zhu directs Nick to prioritize the previously supplied papers and
  establishes a recurring meeting structure. \\
\addlinespace
May 18--20 & Advisor direction & High &
  Professor Zhu connects Nick with Wenbin to clarify and outline the research
  objectives based on the proposal. \\
\addlinespace
June 11 & Meeting logistics & Contextual &
  Confirms the in-person meeting with Wenbin in the Teresa Heinz collaboration
  space. \\
\addlinespace
June 12 & Advisor terminology guidance & High &
  Records that Wenbin requested the phrase ``post-hoc inference'' in the formal
  project description. \\
\addlinespace
June 25--July 3 &
  Student proposal followed by advisor reformulation & Critical &
  Nick proposes a primitive-operation framework; after the June~29 meeting,
  Wenbin drafts a problem setup that ``aligns closer to what we intended.'' \\
\bottomrule
\end{tabularx}
\end{center}

\clearpage
\section{Initial Reading Priorities and Working Structure}
\CriticalSource{Email thread 1 (May 11--14)}

\Meta{Date}{May 11--14, 2026}
\Meta{Participants}{Nick Mino; Professor Woody Zhu}
\Meta{Subject / thread}{Starting the research work effectively}
\Meta{Authority}{Advisor direction responding to student request}
\Meta{Importance}{High}
\Meta{Status}{Foundational background direction; later narrowed by subsequent
guidance}

\paragraph{Why this matters.}
This thread establishes the first concrete working instructions after the
project began. Nick asks which materials and concepts to prioritize, and
Professor Zhu directs him to work through the papers already supplied before
reconnecting. Professor Zhu also indicates that a PhD student will provide
closer guidance as the project develops.

\paragraph{Key contributions.}
\begin{itemize}[leftmargin=*,itemsep=0.25em]
  \item Establishes the initial literature-reading phase.
  \item Creates an expectation of recurring meetings.
  \item Foreshadows the later introduction of Wenbin as the primary technical
        guide.
\end{itemize}

\paragraph{Relationship to the research direction.}
This is background-setting rather than the final research formulation. It
explains why the project initially centered on uncertainty quantification and
the group's existing papers.

\subsection*{Full email thread}
\Email{Thu, May 14, 2026 at 12:45 AM}{Woody Zhu
\texttt{<shixianz@andrew.cmu.edu>}}{%
Sure thing. We will organize meetings on a weekly basis. Sorry for my delayed
response --- I have been traveling lately.

In the meantime, please go through the papers I sent earlier, and let's
reconnect next Wednesday once I return to Pittsburgh.

After you finish the reading, I will connect you with one of our PhD students,
who can provide more guidance as you get further into the project.

Lastly, it is also perfectly fine to mention the research experience on
LinkedIn.

Best,\\
Woody}

\Email{Mon, May 11, 2026 at 1:39 PM}{Nick Mino
\texttt{<nmino@andrew.cmu.edu>}}{%
Hello Professor Zhu,

I wanted to ask what you would recommend I do first to start the research work
effectively.

I have already started reading additional papers and trying to build a
stronger general understanding of uncertainty quantification, but I want to
make sure I am prioritizing the right background material and not moving in an
unproductive direction. Are there specific papers, notes, codebases, or
concepts you think I should focus on first?

I also wanted to ask what the expected structure will be for the research. Will
there be in-person meetings, online meetings, asynchronous materials, or some
combination of those? I am generally available after 12:30~PM on weekdays and
flexible outside of that, so I should be able to adjust to whatever meeting
format works best.

Separately, I wanted to ask how I should describe the role publicly and
professionally. Would it be appropriate for me to make a LinkedIn post
mentioning that I am starting research work, or would you prefer that I wait?
Also, for LinkedIn and my resume, what title would be most accurate:
``Research Assistant,'' ``Undergraduate Research Assistant,'' or something
else?

Best,\\
Nick}

\clearpage
\section{Professor Zhu Connects Nick with Wenbin}
\CriticalSource{Email thread 2 (May 18--20)}

\Meta{Date}{May 18--20, 2026}
\Meta{Participants}{Nick Mino; Professor Woody Zhu; Wenbin Zhou}
\Meta{Subject / thread}{Paper priorities and research-objective clarification}
\Meta{Authority}{Direct advisor instruction}
\Meta{Importance}{High}
\Meta{Status}{Active guidance; establishes Wenbin's role}

\paragraph{Why this matters.}
Professor Zhu explicitly asks Nick to work with Wenbin to clarify and outline
the overall research objectives based on the proposal. This is the formal
handoff from broad reading under Professor Zhu to closer technical development
with Wenbin.

\paragraph{Key contributions.}
\begin{itemize}[leftmargin=*,itemsep=0.25em]
  \item Wenbin is introduced as a third-year PhD student and co-author on
        several relevant papers.
  \item Nick is instructed to clarify and outline the research objectives
        before the next meeting with Professor Zhu.
  \item Nick identifies the specific papers he had been reading: inverse
        conformal risk control, conformalized decision risk assessment,
        Gen-DFL, and hierarchical probabilistic conformal prediction.
\end{itemize}

\paragraph{Relationship to the research direction.}
This thread explains why later messages from Wenbin carry substantial authority
over the technical problem formulation. It also records the literature base
from which the project initially developed.

\subsection*{Full email thread}
\Email{Wed, May 20, 2026 at 11:04 AM}{Woody Zhu
\texttt{<shixianz@andrew.cmu.edu>}}{%
Thank you, Wenbin!

Best,\\
Woody}

\Email{Wed, May 20, 2026 at 7:26 AM}{Wenbin Zhou
\texttt{<wenbinz2@andrew.cmu.edu>}}{%
Thank you, Woody, for connecting us.

Nick has reached out to me in a separate email, and I will chat more with him
offline.

Best,\\
Wenbin}

\Email{Tue, May 19, 2026 at 8:12 PM}{Woody Zhu
\texttt{<shixianz@andrew.cmu.edu>}}{%
Hi Nick,

Sorry for the delayed response. I've been traveling recently and will be fully
back on campus after June~1st. After that, we can resume meeting on a more
regular weekly basis.

In the meantime, I'd like to connect you with Wenbin Zhou, who is a third-year
PhD student in our group and also a co-author on several of these papers.

Please feel free to reach out to Wenbin if you have questions about the papers
or related background material. I would also suggest working with him to
further clarify and outline the overall research objectives based on the
proposal before we meet again. I think that would help make our future
discussions more focused and productive.

Thanks,\\
Woody}

\Email{Mon, May 18, 2026 at 12:07 PM}{Nick Mino
\texttt{<nmino@andrew.cmu.edu>}}{%
Hello Professor Zhu,

Sorry for my late response, and I apologize for not making more progress on
the reading sooner.

I have continued working through the material we discussed. I focused most
closely on your papers on inverse conformal risk control and conformalized
decision risk assessment, along with some related papers.

In your last email, you mentioned the papers you had sent earlier. I may have
missed something, but I could not find a separate list or attachments. I have
mainly been reading the papers we had discussed previously and some of your
related work, including Gen-DFL and Hierarchical Probabilistic Conformal
Prediction for Distributed Energy Resources Adoption.

Could you clarify which specific papers I should focus on before we reconnect?

Best,\\
Nick}

\clearpage
\section{Evidence of the June 11 In-Person Meeting}

\Meta{Date}{June 8--11, 2026}
\Meta{Participants}{Nick Mino; Wenbin Zhou}
\Meta{Subject / thread}{Meeting location in Hamburg Hall}
\Meta{Authority}{Administrative corroboration}
\Meta{Importance}{Contextual but important}
\Meta{Status}{Establishes occurrence and location; not substantive research
guidance}

\paragraph{Why this matters.}
This short thread confirms that Nick and Wenbin met in person on June~11, 2026,
in the Teresa Heinz collaboration space on the B-level of Hamburg Hall. The
substantive research insight from this meeting is preserved separately in
Nick's contemporaneous notes, because no transcript is available.

\paragraph{Key contributions.}
\begin{itemize}[leftmargin=*,itemsep=0.25em]
  \item Confirms the date and physical location of the meeting.
  \item Supports the provenance of the June~11 contemporaneous research note.
  \item Should be cross-referenced by the main research packet.
\end{itemize}

\paragraph{Relationship to the research direction.}
The thread itself contains no technical research content. Its value is
evidentiary: it validates that the undocumented in-person meeting occurred.

\subsection*{Full email thread}
\Email{Thu, Jun 11, 2026 at 3:07 PM}{Wenbin Zhou
\texttt{<wenbinz2@andrew.cmu.edu>}}{%
Hey Nick,

I'm at the Teresa Heinz collab space.

Best,\\
Wenbin}

\Email{Thu, Jun 11, 2026 at 3:05 PM}{Nick Mino
\texttt{<nmino@andrew.cmu.edu>}}{%
I apologize for not asking this sooner, but where in Hamburg b level are we
meeting? I have been in the Teresa Heinz collab space.}

\Email{Tue, Jun 9, 2026 at 11:41 AM}{Nick Mino
\texttt{<nmino@andrew.cmu.edu>}}{%
B-level Hamburg Hall works for me.

Best,\\
Nick}

\Email{Mon, Jun 8, 2026 at 1:57 PM}{Wenbin Zhou
\texttt{<wenbinz2@andrew.cmu.edu>}}{%
Thanks. How about the basement (B-level) of Hamburg Hall?}

\clearpage
\section{Required Terminology: ``Post-Hoc Inference''}
\CriticalSource{Email thread 4 (June 12 terminology)}

\Meta{Date}{June 12, 2026}
\Meta{Participants}{Nick Mino; Wenbin Zhou}
\Meta{Subject / thread}{SURA project description and signed form}
\Meta{Authority}{Advisor terminology guidance embedded in administrative
exchange}
\Meta{Importance}{High}
\Meta{Status}{Active terminology requirement}

\paragraph{Why this matters.}
Although this exchange primarily concerns a signed form, Nick explicitly states
that he incorporated Wenbin's suggestion to add the phrase ``post-hoc
inference'' to the project description. This is direct evidence that the
inferential objective was intended to be part of the formal research framing.

\paragraph{Key contributions.}
\begin{itemize}[leftmargin=*,itemsep=0.25em]
  \item Records Wenbin's requested terminology.
  \item Shows that the project was not merely about generic stability or robust
        optimization.
  \item Supports using post-hoc inference as a defining component of the
        current problem statement.
\end{itemize}

\paragraph{Relationship to the research direction.}
This terminology constraint remains consistent with the later
optimization-stability formulation. It should be treated as an active
requirement unless superseded by explicit advisor guidance.

\subsection*{Full email thread}
\Email{Fri, Jun 12, 2026 at 3:10 PM}{Nick Mino
\texttt{<nmino@andrew.cmu.edu>}}{%
Hi Wenbin,

Thank you, this version works. I appreciate your help.

Best,\\
Nick}

\Email{Fri, Jun 12, 2026 at 3:01 PM}{Wenbin Zhou
\texttt{<wenbinz2@andrew.cmu.edu>}}{%
Here is the PDF version.}

\Email{Fri, Jun 12, 2026 at 2:59 PM}{Wenbin Zhou
\texttt{<wenbinz2@andrew.cmu.edu>}}{%
Let me know if this works. There was probably some formatting issue that caused
the signature to be unable to display.}

\Email{Fri, Jun 12, 2026 at 2:56 PM}{Nick Mino
\texttt{<nmino@andrew.cmu.edu>}}{%
Hi Wenbin,

Thank you for sending this back. I incorporated your suggestion and added
``post-hoc inference'' to the project description.

I noticed that the phone number and date were added, but the additional mentor
signature line at the bottom is still blank. Would you be able to sign that
line in the attached updated version?

Best,\\
Nick}

\clearpage
\section{Primitive-Operation Proposal and Advisor Reformulation}
\CriticalSource{Email thread 5 (June 25--July 3) --- critical}

\Meta{Date}{June 25--July 3, 2026}
\Meta{Participants}{Nick Mino; Professor Woody Zhu; Wenbin Zhou}
\Meta{Subject / thread}{Research approach, June~29 meeting, and revised problem
setup}
\Meta{Authority}{Student proposal followed by advisor-authored reformulation}
\Meta{Importance}{Critical}
\Meta{Status}{July~2 formulation supersedes earlier student wording where
inconsistent}

\paragraph{Why this matters.}
This is the most consequential email chain in the supplied record. Nick
proposes a framework based on stabilized algorithms, primitive operations,
problem structure, and downstream inference goals. After the June~29 meeting,
Wenbin states that he drafted a new problem setup in Overleaf that ``aligns
closer to what we intended'' and authorizes Nick to begin brainstorming
methodologies for that formulation.

\paragraph{Key contributions.}
\begin{itemize}[leftmargin=*,itemsep=0.25em]
  \item Documents Nick's primitive-operation and stabilization-pattern proposal
        before advisor review.
  \item Establishes that the June~29 meeting was used to evaluate the current
        approach and SURA proposal diagram.
  \item Contains Wenbin's explicit statement that his revised setup is closer
        to the advisors' intended problem.
  \item Moves the project from problem clarification into methodology
        development.
\end{itemize}

\paragraph{Relationship to the research direction.}
Nick's June~25 message is evidence of the proposed architecture, not conclusive
evidence of advisor approval. Wenbin's July~2 Overleaf formulation has higher
authority and should control the current problem statement. The attached files
and the actual Overleaf text are therefore essential companion evidence.

\subsection*{Full email thread}
\Email{Fri, Jul 3, 2026 at 2:24 AM}{Nick Mino
\texttt{<nmino@andrew.cmu.edu>}}{%
Hi Wenbin,

Thank you for drafting the problem setup. I will read through it carefully and
start brainstorming possible methodologies to add to the draft. I will let you
know if I have any questions.

For the recurring meeting, I do not think we have one set up yet. I am
available any weekday before 10:30~AM or after 12:30~PM. A weekly meeting would
work well for me.

Best,\\
Nick}

\Email{Thu, Jul 2, 2026 at 12:53 PM}{Wenbin Zhou
\texttt{<wenbinz2@andrew.cmu.edu>}}{%
Hi Nick,

Nice talking to you on Monday, and I am glad to see that we are making good
progress.

I drafted a version of the problem setup in the Overleaf project, which aligns
closer to what we intended. Please feel free to give it a read and let me know
if you have any questions. If not, you may begin brainstorming the
methodologies we could use to approach this problem and write them down in the
draft. We can discuss them next time we meet.

Woody --- Please see attached for the two files that Nick presented on Monday.

By the way, do we have a weekly recurring meeting with Nick set up? I don't see
any on my calendar.

Best,\\
Wenbin}

\Email{Fri, Jun 26, 2026 at 12:56 PM}{Nick Mino
\texttt{<nmino@andrew.cmu.edu>}}{%
That works for me, thank you.}

\Email{Fri, Jun 26, 2026 at 11:56 AM}{Wenbin Zhou
\texttt{<wenbinz2@andrew.cmu.edu>}}{%
Works for me as well.}

\Email{Fri, Jun 26, 2026 at 10:45 AM}{Woody Zhu
\texttt{<shixianz@andrew.cmu.edu>}}{%
Sorry I realized I have a meeting between 4--5. How about 5--5:30?

Woody}

\Email{Fri, Jun 26, 2026 at 9:04 AM}{Wenbin Zhou
\texttt{<wenbinz2@andrew.cmu.edu>}}{%
Works for me}

\Email{Fri, Jun 26, 2026 at 8:02 AM}{Woody Zhu
\texttt{<shixianz@andrew.cmu.edu>}}{%
What about 4--4:30?

Woody}

\Email{Fri, Jun 26, 2026 at 6:37 AM}{Wenbin Zhou
\texttt{<wenbinz2@andrew.cmu.edu>}}{%
Hi Woody,

12:30--1:00 is a bit tough for me. Would any time before 12 or after 1 work?

Best,\\
Wenbin}

\Email{Thu, Jun 25, 2026 at 10:09 PM}{Woody Zhu
\texttt{<shixianz@andrew.cmu.edu>}}{%
Hi Nick,

Thanks for reaching out. How about next Monday 12:30--1?

Let me know if this works for you and Wenbin.

Best,\\
Woody}

\Email{Thu, Jun 25, 2026 at 4:19 PM}{Nick Mino
\texttt{<nmino@andrew.cmu.edu>}}{%
Hello Professor Zhu and Wenbin,

I hope you are doing well. I wanted to ask whether we could schedule a meeting
soon to discuss my current research approach and the SURA Week~3 proposal
diagram.

My current approach is to focus on algorithmic stability for post hoc
selection/inference by studying concrete examples of selection algorithms that
have been `stabilized'. I have been gathering these examples so that I can
compare the structure of the original selection algorithm with the stabilized
version, identify the primitive operations involved, and look for recurring
patterns in how stability is introduced. I am in the process of converting the
examples I found into LaTeX.

The goal is to begin work on a framework that analyzes a selection algorithm's
operations, the relevant data type or problem structure, and the downstream
inference or evaluation goal to determine an appropriate formulation for adding
algorithmic stability.

I wanted to check whether this is the right approach to continue pursuing.

For the Week~3 assignment, I also need to ask a few questions about the
research question, novelty, significance, methods, and expected outcomes, so
the next meeting would be useful for confirming those pieces.

I also wanted to ask what schedule would make sense for regular meetings going
forward. Would you prefer weekly or biweekly meetings, and is there a recurring
time that generally works for both of you?

I am available any day before 10:30~AM or after 12:30~PM, except Saturdays.

Best,\\
Nick}

\clearpage
\section*{Concluding Synthesis: Evolution of Advisor Guidance}
\begin{enumerate}[leftmargin=*,itemsep=0.6em]
  \item \textbf{Initial background phase.}
        Professor Zhu first directed Nick to study the group's
        uncertainty-quantification and decision-focused papers. At this stage,
        the technical problem remained broad.

  \item \textbf{Technical mentorship and objective clarification.}
        Professor Zhu then connected Nick with Wenbin and explicitly asked them
        to clarify and outline the research objectives based on the proposal.
        This established Wenbin as the primary source of detailed technical
        guidance.

  \item \textbf{Post-hoc inference becomes explicit.}
        The June~12 exchange records Wenbin's suggestion that the formal
        project description include the phrase ``post-hoc inference.'' This
        indicates that downstream inferential validity is a defining objective,
        not an optional application.

  \item \textbf{Student-proposed primitive framework.}
        By June~25, Nick proposed studying concrete stabilized selection
        algorithms, decomposing them into primitive operations, and identifying
        recurring stabilization patterns. This is the clearest email statement
        of the architecture Nick brought to the advisors.

  \item \textbf{Advisor reformulation controls.}
        After the June~29 meeting, Wenbin drafted a new problem setup that he
        said aligned more closely with what the advisors intended. That
        formulation supersedes earlier student wording where the two differ and
        should be the primary source for the current mathematical problem.

  \item \textbf{Methodology phase begins.}
        Wenbin then instructed Nick to begin brainstorming and writing
        methodologies for the revised setup. The project had therefore moved
        from broad reading and problem identification into the development of
        candidate methods.
\end{enumerate}

\Takeaway{The project concerns algorithmic stability for optimization-based
selection in support of post-hoc inference, with the final formulation governed
by the advisor-authored July~2 problem setup. The primitive-operation framework
is an important proposed methodology, but its status must be evaluated against
that advisor-authored formulation and the June~29 meeting transcript.}

\section*{Missing Companion Evidence}
\begin{itemize}[leftmargin=*,itemsep=0.35em]
  \item The actual July~2 Overleaf problem setup.
  \item The two files attached to Wenbin's July~2 email.
  \item The June~29 meeting transcript.
  \item The June~11 contemporaneous notes describing the optimization-specific
        stability intuition.
  \item Any earlier June~5 email in which Wenbin directly connected algorithmic
        stability to replacing data splitting in CREME, if available.
\end{itemize}

\begin{center}
\textit{End of record.}
\end{center}

\end{document}
